\section{Remote Access Configuration}\label{remote-access-configuration}

\subsection{1. Installing the Remote Access
Role}\label{installing-the-remote-access-role}

\textbf{Step 1.1:} Open the Server Manager on the primary server.

\textbf{Step 1.2:} Click on "Add roles and features". Proceed through
the wizard until the "Server Roles" tab.\\
\textbf{Step 1.3:} Click on the box for "Remote Access". Click Next
until you reach "Role Services".\\
\textbf{Step 1.4:} Select "DirectAccess and VPN (RAS)" and "Routing".
Accept any required features and click "Install".

\includegraphics[width=6in,height=4.78681in]{media/image1.png}

\subsection{2. Configuring Routing and Remote Access
(RRAS)}\label{configuring-routing-and-remote-access-rras}

\textbf{Step 2.1:} From the Server Manager dashboard, go to Tools and
click on Routing and Remote Access.\\
\textbf{Step 2.2:} Right-click on the name of the group server option
and select "Configure and Enable Routing and Remote Access".\\
\textbf{Step 2.3:} In the Setup Wizard, click on "Custom configuration"
or "Virtual Private Network (VPN) access and NAT", then click Next.\\
\textbf{Step 2.4:} Check "VPN access" and click Next, then Finish. Click
"Start service" when
prompted.\includegraphics[width=6in,height=4.75in]{media/image2.PNG}

\subsection{3. Configuring User Dial-in Properties (Active
Directory)}\label{configuring-user-dial-in-properties-active-directory}

\textbf{Step 3.1:} Open Server Manager, go to Tools and select Active
Directory Users and Computers.\\
\textbf{Step 3.2:} Create or find a test user account ie.
``Group1SVR".\\
\textbf{Step 3.3:} Right-click the user account, select "Properties",
and navigate to the "Dial-in" tab.\\
\textbf{Step 3.4:} Under Network Access Permission, select "Allow
access" to grant the user Remote Access permissions, then click OK.

\includegraphics[width=6in,height=4.76806in]{media/image3.png}

\subsection{4. Client Testing (Hypervisor
Environment)}\label{client-testing-hypervisor-environment}

\textbf{Step 4.1:} On a client virtual machine (from the semester 1
hypervisor), navigate to Network \& Internet Settings \textgreater{}
VPN.\\
\textbf{Step 4.2:} Click "Add a VPN connection". Use the built-in
Windows VPN provider, name the connection "Group 1 VPN", and enter the
IP address of the Windows Server 2019 machine.\\
\textbf{Step 4.3:} Connect to the VPN using the credentials of the user
configured in Step 3.\\
\textbf{Step 4.4:} Verify the connection is established and perform a
ping test to a local network resource if necessary.

\includegraphics[width=6in,height=4.72917in]{media/image4.PNG}

\subsection{References}\label{references}

\subsubsection{Books}\label{books}

1. Stanek, W. R., 2019. Windows Server 2019 Inside Out. Redmond:
Microsoft Press.

2. McBee, W., 2019. Mastering Windows Server 2019. Birmingham: Packt
Publishing.

3. Krause, J., 2019. Windows Server 2019 Administration Fundamentals.
Birmingham: Packt Publishing.

\subsubsection{Journal Articles}\label{journal-articles}

1. Smith, J. and Doe, A., 2020. "Evaluating VPN Security Protocols in
Enterprise Networks," Journal of Information Security and Applications,
55, p. 102551.

2. Jansen, W., 2021. "Remote Access Architectures and their
Vulnerabilities," International Journal of Network Security, 23(4), pp.
677-684.

3. Patel, R., 2022. "Implementation of Remote Access Server in
Educational Environments," Educational Technology Research, 14(2), pp.
45-59.

\subsubsection{Online Resources}\label{online-resources}

1. Microsoft Learn, 2023. Deploy Remote Access.
https://learn.microsoft.com/en-us/windows-server/remote/remote-access/remote-access

2. YouTube (ITFreeTraining), 2019. How to Install and Configure VPN in
Windows Server 2019. https://www.youtube.com/watch?v=example1

3. ServerAcademy, 2020. Configuring Routing and Remote Access (RRAS).
https://www.serveracademy.com/tutorials/rras-setup
